% ============================================================
%  Abstract
% ============================================================
\chapter*{Abstract}
\addcontentsline{toc}{chapter}{Abstract}

Decentralized multi-robot task allocation (MRTA) methods such as
CBBA~\cite{choi2009consensus} and its successor
GCBBA~\cite{kha2025enhancing} assume fully connected or at least static
communication graphs.
These assumptions rarely hold in deployments where range-limited radios,
obstacles, and robot mobility create dynamic, intermittently disconnected
topologies.
This thesis presents the \textbf{Local Consensus-Based Bundle Algorithm (LCBA)},
which extends GCBBA by replacing its global consensus requirement with
component-local consensus.
LCBA relaxes three key assumptions of its predecessors:
(1)~tasks need not be known \emph{a priori}; LCBA handles dynamic,
steady-state arrivals via event-driven replanning;
(2)~the communication graph may be arbitrarily disconnected; agents restrict
consensus to their current connected component and resolve conflicts upon
reconnection; and
(3)~the communication topology may change at every timestep.

LCBA is integrated with three interchangeable path planning strategies
(Cooperative A*, Priority-Based Search, and Rolling Horizon Collision
Resolution) forming a complete multi-robot coordination pipeline with a
modular, planner-agnostic allocation layer.
The system is evaluated across six gridworld environments of increasing scale
(30$\times$30 with 6 agents, 44$\times$28 with 12 agents, 36$\times$36 with
20 agents, 60$\times$60 with 18 agents, 72$\times$72 with 80 agents, and
161$\times$63 with 200 agents) against three baselines: CBBA, the centralized
Sequential Greedy Algorithm~(SGA), and the state-of-the-art decentralized
method DMCHBA~\cite{samiei2024dmchba}.

Experiments are conducted in two modes.
In batch mode (fixed task set, makespan objective), LCBA achieves
non-zero task completion at all six communication ranges tested, including
settings where CBBA and SGA achieve zero completion rate.
In steady-state mode (continuous task arrivals, throughput objective),
LCBA's throughput is flat across the full communication spectrum
(from fully disconnected to fully connected) while all baselines degrade
sharply below the connectivity threshold.
On connected graphs, LCBA achieves a sub-optimality ratio of
\textbf{1.0--1.3} relative to SGA, within the theoretical
50\%-optimality bound inherited from GCBBA.
The path planner comparison reveals that CA* offers the best
computation-throughput trade-off for small fleets, while RHCR is preferred
under heavy load or for larger agent counts.

The results establish LCBA as the most robust decentralized allocation
method across the full communication spectrum, with implications for
warehouse automation, disaster response, and other domains where
communication reliability cannot be guaranteed.
